\chapter{Vulnerabilidade, crise e suspensão, 1911--1914}

\section{O auge e seus problemas}

A reforma de dezembro de 1910 ampliou a capacidade da Caixa e preservou a estabilidade a uma taxa mais alta. Em 1911, o regime podia ser apresentado como realização duradoura: depósitos expressivos, câmbio fixo, expansão da circulação e acesso a capitais externos. O auge, porém, trouxe novas questões sobre onde o ouro era mantido, quem tinha autoridade sobre ele e que obrigações ligavam governo, depositantes e agentes financeiros.

O episódio mais revelador ocorreu quando o governo transferiu depósitos de ouro da Caixa aos Rothschild, com remuneração anual de 2,5 por cento destinada ao fundo de garantia. Os quatro jornais trataram do ato. A controvérsia não dizia respeito diretamente a aumentar ou reduzir a taxa, mas à natureza jurídica do depósito. As moedas haviam sido entregues por particulares em troca de notas, e a lei determinava que o lastro fosse conservado para o resgate. Investi-lo no exterior podia ser descrito como administração prudente de uma reserva ou como disposição indevida de patrimônio confiado ao Estado.

\fonteaconferir{Reunir as matérias dos quatro jornais sobre a transferência de 1911, conferir data, valor, instrumento jurídico, taxa de 2,5 por cento e destino dos juros. Confrontar o ato com o art. 1º, parágrafo 2º, do Decreto n. 1.575.}

O \textit{Correio da Manhã} atacou a medida pelo direito de propriedade, sustentando que o governo não podia usar livremente ouro entregue por particulares. Outros textos enfatizaram segurança, rendimento e confiança na casa bancária. O episódio oferece uma janela para quatro dimensões que o eixo entre valorização e expansão não capta: custódia, soberania, confiança e poder financeiro.

Ele também permite qualificar a presença dos Rothschild no debate, e essa presença foi medida. Nas 8.331 páginas do corpus que mencionam a Caixa de Conversão, 170 nomeiam um Rothschild, enquanto 456 designam o credor externo por cargo ou perífrase, como \textit{agentes financeiros}, \textit{banqueiros em Londres} ou expressões equivalentes. Das 456, um subconjunto de 214 usa perífrase de confiança alta. Apenas 38 páginas fazem as duas coisas, e 418 designam o credor sem nunca nomeá-lo. A imprensa torna esse ator presente e, ao mesmo tempo, majoritariamente anônimo.

A medida tem um limite que precisa acompanhá-la. O ruído de OCR atinge nome próprio com mais força do que expressão comum, porque um sobrenome estrangeiro não se recupera por contexto quando uma letra falha. A contagem de perífrases é, por isso, mais confiável que a de nomes, e a razão entre as duas deve ser lida como piso da desproporção, não como sua medida exata. Nos volumes parlamentares sobre a Caixa, os banqueiros estrangeiros quase não aparecem pelo nome. A diferença pode refletir regimes de fala: o jornal acompanha operações e telegramas cotidianos, enquanto o plenário discute competências e princípios em linguagem nacional.

\interpretacaoprovisoria{A transferência de 1911 não será apresentada como causa da crise macroeconômica posterior. Sua importância é institucional e discursiva: ela torna visíveis as condições de confiança do regime e a ambiguidade entre ouro como lastro intocável e reserva administrada por governo e agentes financeiros.}

\section{Estabilidade sob observação em 1912}

Em meados de janeiro de 1912, retiradas de aproximadamente três mil contos provocaram alarme na praça. A \textit{Gazeta de Notícias} procurou reduzir a gravidade do movimento, comparando-o ao total dos depósitos e insistindo que a Caixa era instituição de emissão, não banco de depósitos. No mesmo editorial, o jornal declarou sua responsabilidade no apoio prestado à instituição e recuperou a oposição à política de Bulhões. Essa autodescrição oferece evidência rara de apropriação editorial explícita.

\fonteaconferir{O editorial é de 17 de janeiro de 1912, e não de 15, como o rascunho registrava. Objeto \texttt{per103730\_1912\_00017}, página 1. A data foi fixada pelos mastheads vizinhos, \texttt{00016} traz terça-feira 16 e \texttt{00018} traz 18 de janeiro, e confere com o calendário, em que 17 de janeiro de 1912 caiu numa quarta-feira. O objeto de 15 de janeiro traz apenas o boletim de rotina da Caixa. Falta conferir na imagem o valor das retiradas e a proporção em relação às reservas. A extração tratada apresenta sangria de coluna no terço final da peça, que se interleava com uma entrevista sobre a revolução paraguaia.}

O episódio sugere que saídas de ouro ainda podiam ser tratadas como rotina reversível. A crença na estabilidade não exigia negar qualquer retirada, mas enquadrá-la como movimento pequeno diante de reservas suficientes. Importa observar quando esse vocabulário de normalização perde força e é substituído por termos de corrida, escassez, defesa ou emergência.

A partir de 1912, outros sinais se acumularam. A situação fiscal se deteriorava, e o governo começou a encontrar dificuldades para contratar novos empréstimos. Credores passaram a demonstrar preocupação com déficits, compromissos estaduais e capacidade de pagamento. O debate sobre empréstimos externos dos estados opôs autonomia federativa à necessidade de coordenação nacional para preservar o câmbio e o crédito do país.

\fonteaconferir{Localizar as peças de 1912 e 1913 sobre empréstimos estaduais, inclusive o argumento de 5 de novembro de 1913 segundo o qual operações sem autorização federal assustavam capitalistas europeus.}

\section{A reversão antes da guerra}

A crise da Caixa não começou em agosto de 1914. Seus mecanismos se formaram antes da guerra, sobretudo em 1912 e 1913. A explicação exige combinar vários choques em vez de atribuir causalidade exclusiva às Guerras Balcânicas.

O primeiro componente foi exportador. A borracha brasileira perdeu competitividade diante da produção asiática, e seu preço caiu acentuadamente depois do auge de 1910. Como o produto havia se tornado uma fonte importante de divisas, a queda reduziu receitas externas e atingiu regiões e governos dependentes da atividade. O café também enfrentou problema político nos Estados Unidos. A ação antitruste contra a valorização levou ao embargo de estoques mantidos em Nova York, e o comitê de credores decidiu vendê-los em janeiro de 1913, pressionando preços e alterando o controle do programa \cite{fritsch1990,francolago2011}.

O segundo componente foi financeiro. As Guerras Balcânicas elevaram a incerteza e restringiram a oferta europeia de capitais. Para uma economia que financiava investimento, gasto público e estoques por empréstimos externos, a interrupção não era apenas perda de recurso adicional. Ela removia a entrada que compensava importações e sustentava reservas. A piora fiscal e a dificuldade de novas operações reforçavam a cautela dos credores.

O terceiro componente foi monetário. O déficit externo e a retirada de ouro reduziam as notas da Caixa em circulação. Como os mercados domésticos de capital eram estreitos e a esterilização limitada, a contração do lastro se transmitia ao crédito e à atividade. A recessão, por sua vez, enfraquecia receitas públicas e bancos, ampliando a demanda por liquidez. Esse é o mecanismo pelo qual o padrão-ouro na margem podia intensificar uma reversão originada nas exportações e nos capitais \cite{oliveirasilva2001,fritschfranco1992}.

Segundo Oliveira e Silva \cite{oliveirasilva2001}, a queda de café e borracha, a redução dos capitais depois das Guerras Balcânicas, a perda de reservas e a contração monetária formaram o núcleo da reversão. Fritsch \cite{fritsch1990} oferece uma narrativa mais abrangente, incluindo a deterioração fiscal, o episódio antitruste, as dificuldades de empréstimo e a recessão já instalada antes da Primeira Guerra. As explicações são complementares em parte, mas diferem no peso e no encadeamento dos fatores. A redação final deverá confrontar séries de reservas, circulação, exportações, preços e empréstimos para datar cada inflexão.

\revisarbibliografia{Verificar a edição de Fritsch usada para a cronologia de 1912 a 1914 e incorporar literatura específica sobre o choque da borracha, a ação antitruste e a crise financeira brasileira de 1913.}

A ausência da \textit{Gazeta de Notícias} em 1913 limita a comparação entre os quatro jornais justamente no ano da reversão. O silêncio dessa célula é documental, não editorial. As posições dos outros três periódicos deverão ser comparadas sem imputar à \textit{Gazeta} a continuidade de seu editorial de 1912.

\section{Quando os compromissos se tornam incompatíveis}

A expressão incompatibilidade de compromissos é uma síntese analítica desta dissertação. Ela não pretende atribuir à literatura uma fórmula literal. Designa uma situação em que objetivos anteriormente conciliados pela abundância externa passaram a exigir decisões conflitantes.

A conversibilidade imediata exigia entregar ouro a quem apresentasse notas da Caixa. A preservação do lastro recomendava limitar a saída para conservar a reserva restante. A defesa da taxa de 16 implicava aceitar contração de notas e crédito enquanto o balanço de pagamentos permanecesse negativo. A liquidez bancária exigia fornecer meios de pagamento num momento de retirada e incerteza. A atividade econômica e o emprego sofriam com a contração. As finanças públicas perdiam receita e precisavam sustentar bancos, governos e obrigações. O serviço da dívida externa continuava exigindo divisas e reputação contratual.

Em conjuntura favorável, entrada de ouro, expansão, estabilidade e confiança podiam se reforçar. Na reversão, preservar cada compromisso integralmente tornava mais difícil cumprir os demais. Emitir sem ouro aliviaria bancos e Tesouro, mas romperia a regra da Caixa. Manter o troco até o fim protegeria o contrato individual, mas poderia esgotar o lastro. Suspender preservaria a reserva, mas negaria temporariamente a conversão prometida. Defender a paridade por contração poderia aprofundar a recessão e reduzir a própria capacidade fiscal de honrar compromissos externos.

Essa colisão ajuda a formular a questão da crença monetária. A crise não precisa transformar defensores da estabilidade em adversários do padrão-ouro. Ela pode alterar a percepção sobre viabilidade, prioridade e custo. Um ator podia continuar considerando desejável a moeda estável e, ainda assim, julgar inviável a conversão imediata. Podia defender suspensão temporária para restaurar a paridade depois. Podia rejeitar a suspensão por honra contratual ou rejeitá-la porque preferia deixar o ouro se esgotar e encerrar o câmbio fixo. A posição sobre uma medida emergencial não identifica sozinha o regime desejado.

\section{Agosto de 1914}

A eclosão da guerra fechou ou desorganizou mercados financeiros, interrompeu remessas e ampliou a procura por liquidez. O governo decretou feriados bancários no começo de agosto e adotou medidas para conter corridas e dar tempo às autoridades. Em 15 de agosto, o Decreto n. 2.862 suspendeu o troco das notas da Caixa por ouro \cite{brasil1914suspensao}. A medida marcou o fim operacional da conversibilidade no recorte desta dissertação, embora a instituição e suas obrigações continuassem a produzir efeitos jurídicos e políticos.

Suspender o troco preservava o ouro restante, mas alterava o significado das notas. A reserva deixava de estar imediatamente disponível aos portadores e podia ser tratada como recurso estratégico diante da guerra. A medida era emergencial e inicialmente temporária, o que deixava aberta a expectativa de retorno. Prorrogações posteriores, contudo, ampliaram a distância entre promessa original e operação efetiva.

Em 24 de agosto, o Decreto n. 11.091 autorizou emissão de 150 mil contos, dos quais 100 mil seriam destinados a empréstimos aos bancos \cite{brasil1914emissao}. A emissão de socorro respondia à escassez monetária e à ameaça de crise bancária. Ela não deve ser lida apenas como escolha doutrinária. Era intervenção de emergência numa economia que já havia contraído antes da guerra. Seu tamanho e efeito líquido precisam ser comparados à redução anterior da circulação.

O governo também negociou um novo \textit{funding loan}, concluído no segundo semestre, para reorganizar compromissos externos. Assim como em 1898, o refinanciamento ligou crise cambial, finanças públicas e credores. O paralelo é importante, mas as conjunturas eram diferentes. Em 1914, a ruptura atingia o próprio sistema internacional que sustentava o padrão-ouro, e não apenas a reputação ou o balanço brasileiro \cite{fritsch1990,francolago2011}.

\fonteaconferir{Conferir decretos de feriado e moratória, prorrogações da suspensão, termos e data do \textit{funding loan} de 1914 e cronologia diária das medidas de agosto.}

\section{O debate sobre a suspensão}

As peças já catalogadas identificam três objetos: suspender o troco, taxar a exportação de ouro como alternativa e definir a responsabilidade pelo excesso de notas sobre o depósito. A primeira controvérsia foi aberta; as outras ainda estão pouco documentadas.

Na sessão de 12 de dezembro, reproduzida pela \textit{Gazeta de Notícias}, Martim Francisco e Serzedello Corrêa votaram contra a prorrogação da suspensão. O resultado formal era igual, mas as razões eram opostas. Martim Francisco invocou a palavra do Estado e a obrigação do depositário de devolver o que recebera. Serzedello considerava a Caixa um obstáculo e queria que funcionasse até o esgotamento do ouro, quando a experiência terminaria. Um defendia a força do contrato; o outro recusava preservar artificialmente o regime.

\fonteaconferir{Conferir a \textit{Gazeta de Notícias} de dezembro de 1914 e os anais da sessão de 12 de dezembro, preservando as falas de Martim Francisco e Serzedello Corrêa como vozes parlamentares reproduzidas, não editoriais do jornal.}

O caso demonstra por que a variável favorável ou contrário à suspensão é insuficiente. Para interpretar a posição, é preciso registrar o veredito retrospectivo sobre a Caixa, o regime desejado depois da crise, a concepção de depósito e a natureza atribuída à emissão de socorro.

Carlos Peixoto Filho teria preferido imposto sobre a exportação de ouro amoedado à suspensão do troco, invocando precedente do início da Caixa. Martim Francisco também chamou atenção para notas emitidas além do depósito, em montante próximo de 19.890 contos, e para a responsabilidade correspondente. Esses argumentos deslocavam a controvérsia da mera preservação do metal para a legalidade da circulação e as obrigações do Estado.

\fonteaconferir{Confirmar o valor do excedente, a origem contábil e o precedente do imposto sobre exportação de ouro. Localizar o parecer de Carlos Peixoto Filho e as emendas votadas.}

As posições editoriais dos quatro jornais em 1914 ainda não podem ser sintetizadas com a mesma segurança de 1906. Há apenas duas peças classificadas provisoriamente como editoriais substantivos na amostra da fase. Notícias parlamentares e artigos são indispensáveis, mas exigem atribuição cuidadosa de voz.

\lacuna{Ampliar a leitura comparada de julho a dezembro de 1914 nos quatro jornais, com atenção a títulos, comentários próprios, republicações, seções livres e continuidade após as primeiras medidas de agosto.}

\section{O que mudou na crença na estabilidade}

O encerramento do recorte não deve procurar uma conversão coletiva contra o metalismo. A pergunta mais produtiva é como a crise desagregou elementos antes apresentados como um único programa. Entre 1906 e 1910, estabilidade, entrada de ouro, expansão da circulação, confiança e crescimento podiam ser narrados como resultados convergentes. Entre 1912 e 1914, desejabilidade, viabilidade e custo se separaram.

A crença na estabilidade podia sobreviver como objetivo de longo prazo, mesmo quando a taxa de 16 se tornava indefensável. A crença na conversibilidade podia sobreviver como obrigação moral, ainda que a guerra tornasse seu cumprimento imediato impraticável. A suspensão podia ser interpretada como exceção destinada a proteger o retorno futuro ou como demonstração de que o arranjo era vulnerável. O sentido histórico da crise está nessa abertura de alternativas.

\interpretacaoprovisoria{A hipótese para a leitura final é que 1914 não destruiu de uma só vez o prestígio da estabilidade cambial. A crise revelou que estabilidade, conversibilidade e paridade específica não eram sinônimos, e obrigou os participantes a ordenar compromissos que antes podiam defender em conjunto. A mudança de crença deve ser observada nessa diferenciação, no horizonte de retorno e no regime proposto para depois da emergência.}

Essa formulação também permite conectar o caso brasileiro à crise do padrão-ouro internacional sem dissolver a história nacional numa explicação externa. Choques de exportação, política fiscal, empréstimos estaduais, decisões sobre estoques e estrutura bancária prepararam a vulnerabilidade. As Guerras Balcânicas e a Primeira Guerra atingiram uma economia já em reversão. O regime transmitiu esses choques ao crédito e à circulação, e a imprensa lhes deu significado político.
